\documentclass[border=5pt]{standalone}
\usepackage[utf8]{inputenc}
\usepackage[T1]{fontenc}
\usepackage[spanish]{babel}
\usepackage{tikz}
\usetikzlibrary{arrows.meta, positioning, calc, fit, backgrounds, shadows}
\usepackage{amsmath}

\begin{document}
	

	
	% ============================================================
	% SEGUNDO TIKZPICTURE (MODIFICADO CON ESTILO UNIFICADO)
	% ============================================================
	\tikzset{
		estiloAzul/.style={
			rectangle, rounded corners=4pt,
			draw=blue!60!black, fill=blue!8,
			text width=4.5cm, align=center,
			minimum height=1.2cm, font=\small,
			drop shadow
		},
		estiloNaranja/.style={
			rectangle, rounded corners=4pt,
			draw=orange!70!black, fill=orange!8,
			text width=4.5cm, align=center,
			minimum height=1.2cm, font=\small,
			drop shadow
		},
		estiloTeal/.style={
			rectangle, rounded corners=4pt,
			draw=teal!60!black, fill=teal!10,
			text width=4.5cm, align=center,
			minimum height=1.2cm, font=\small,
			drop shadow
		},
		estiloVerde/.style={
			rectangle, rounded corners=4pt,
			draw=yellow!50!black, fill=yellow!15,
			text width=4.5cm, align=center,
			minimum height=1.2cm, font=\small,
			drop shadow
		},
		estiloMorado/.style={
			rectangle, rounded corners=4pt,
			draw=purple!60!black, fill=purple!8,
			text width=4.5cm, align=center,
			minimum height=1.2cm, font=\small,
			drop shadow
		},
		estiloFinal/.style={
			rectangle, rounded corners=4pt,
			draw=green!50!black, fill=green!12,
			text width=4.5cm, align=center,
			minimum height=1.2cm, font=\small,
			drop shadow
		},
		flechaModerna/.style={-{Stealth[length=2.5mm]}, thick, black}
	}
	
	\vspace{1cm}
	\centering
	
	\scalebox{0.70}{
		\begin{tikzpicture}[node distance=2.2cm and 1.5cm]
			
			% Fila superior
			\node (start) [estiloAzul] {1. Selección de datos \\ Obtención de estructuras de ésteres en PubChem.};
			
			% Parámetros
			\node (p1) [estiloNaranja, below=of start] 
			{2. Estandarización topológica \\ Reordenamiento atómico según el patron O1, C2, O3, R1 y R2.};
			
			\node (p2) [estiloTeal, right=of p1] 
			{3. Generación de matrices Z \\ Conversión de coordenadas cartesianas a coordenadas internas.};
			
			% Cálculo
			\node (calc) [estiloVerde, below=of p2] 
			{4. Cálculos mecanocuánticos \\ Optimización de estructuras $\omega$B97X-D/cc-pVDZ, CPCM.};
			
			% Aplicación
			\node (aplica) [estiloMorado, left=of calc] 
			{5. Análisis de frecuencias \\ Verificación de un mínimo en la superficie de energía potencial.};
			
			\node (end) [estiloFinal, left=of aplica] {6. Obtención de función extendida (WFX).};
			
			% Flechas
			\draw [flechaModerna] (start) -- (p1);
			\draw [flechaModerna] (p1) -- (p2);
			\draw [flechaModerna] (p2) -- (calc);
			\draw [flechaModerna] (calc) -- (aplica);
			\draw [flechaModerna] (aplica) -- (end);
			
		\end{tikzpicture}
	}
	
\end{document}